The work followed an iterative, data-first engineering process rather than a fixed up-front design. Each milestone began by inspecting real recorded data (topic schemas, physical units, coordinate frames and clock domains) before any code was written, and ended with a verification step against that same data once a node or fix was in place. This loop of characterizing the raw data, implementing against it, and then re-verifying the result on that same data was applied consistently across the project: position and orientation corrections on the pointcloud, altimeter data cleaning and the disparity-backend search were all driven by first characterizing the raw sensor output and only then developing the code, rather than relying only on vendor documentation or conventions.

Progress was tracked through screenshots, plots and short video clips of the pipeline's output, which were shared with the supervisors for feedback. The project was developed in a single Git repository, with a main branch for stable code and feature branches for new nodes or fixes. Each commit was accompanied by a descriptive message, and the commit history was used to track progress and decisions made during development.

The technology stack was: ROS~2 Humble with the Zenoh middleware implementation, Python (\texttt{rclpy}), OpenCV (classical and CUDA-accelerated stereo matching), PyTorch and ONNX Runtime (learned stereo-matching backends), and Git~\cite{github} for version control. No dedicated project-management tooling was used; given the single-person, single-supervisor nature of the internship, coordination with supervisors was handled through periodic meetings rather than a formal sprint/ticket process.
